\chapter{Discussion}
\label{chap:discussion}

This chapter discusses the findings derived from examining the Ray-Ban Meta glasses. The examination comprised an overview of the device's hardware and software (Chapter \ref{chap:deviceOverview}), an assessment of the companion app (Chapter \ref{chap:appAnalysis}), analyses of network and Bluetooth traffic (Chapters \ref{chap:networkAnalysis} and \ref{chap:bluetoothAnalysis}), and an overview of the Decryption attempts and analysis (Chapter \ref{chap:decryption}). By combining these perspectives, a clearer image of the privacy implications of the glasses and the concerns of data collection and retention, user control, and bystander consent can be painted. The discussion is organized around key topics, metadata exposure and user activity inference, companion app dependencies, trade-offs between functionalities and privacy, and policy-oriented reflections, including bystander privacy.

\subsection*{Metadata Exposure and User Activity Inference}
The network traffic analysis reveals that, although end-to-end encryption was used for the traffic, mainly via TLS and QUIC protocols, users are still vulnerable through metadata exposure, which allows for inference of application usage and user activities. 

Large packet transfers (more than $10^5$ bytes when grouped in ten-second intervals) usually coincide with interactions such as media transfer or issuing voice commands. Media transfer has the largest traffic volume spike compared to all other interactions. Media transfer also has a distinct sequence of packets that starts with four ICMPv6 packets, which appeared exclusively during media transfer. Other voice commands also have different sequences of packets. For instance, sending a message to a person involves a back-and-forth exchange of around six to eight QUIC packets with payload, followed by roughly four packets of empty QUIC packets, whereas asking for the weather typically would have roughly six consecutive QUIC packets that contain payload followed by four empty ones, which was then continued with around eight QUIC packets with payload. These traffic patterns allow an adversary to intercept the communication to reasonably deduce user actions without the need to decrypt the content. 

A similar risk exists in Bluetooth communications as the glasses rely on Bluetooth for pairing, periodic advertisements, and device management. An adversary capable of sniffing BLE advertisements or tracking the device's Lower Address Part (LAP) in Classic Bluetooth would have a stronger basis for user profiling or inferring specific user interactions based on traffic volume. BLE's use of advertising data, including a Company ID (0x01ab for Meta), simplifies the identification of the glasses and correlates transmission to specific features, such as pairing events or periodic advertisements. On the other hand, Classic Bluetooth captures showed traffic bursts that revealed active usage windows, with most traffic volume during music playback or voice-based audio streaming, including asking the AI to read received messages. The glasses with the potential ability to locally process messages received could be vulnerable if the glasses are stolen or compromised. Furthermore, the lack of user awareness about what data is stored and for how long it is stored raises concerns about transparency and bystander privacy.

Most network transmissions were well-protected using TLS or QUIC protocols. However, the partial success in bypassing certificate pinning revealed additional insights into the data flow. For instance, a decrypted response from ar.graph.meta.com revealed that the app actively transmits details such as device IDs and session states to the server. These identifiers allow Meta's servers to monitor app activity and, when combined with other logs, could be aggregated to infer user behavior. Additionally, the repeated requests to lookaside.facebook.com demonstrated how the app fetches image resources. It is harmless for official product images but raises concerns about how images captured by the user may be processed or stored. These findings reinforce the notion that even with encryption, determined adversaries can extract not only metadata but also content-level information, potentially revealing user activities and exposing bystander information. Although the stronger certificate pinning prevented the decryption for the other endpoints, it leaves the question of the exact nature of the data being collected and retained within the more protected transmissions.



\subsection*{Companion App and Data Retention Concerns}
The static analysis of the companion app for the Ray-Ban Meta glasses showed that it does not request excessive permissions nor exhibit excessive reliance on external third-party services. By targeting specific permissions, an adversary could obtain location coordinates, social relationships, and behavioral patterns. This information could not only be used by adversaries but also aggregated into Meta's data collection, allowing for extensive user profiling that extends beyond the app itself.

Although no explicit references to dedicated advertising SDKs were detected, the heavy usage of domains like graph.facebook.com or spclient-based Spotify APIs and the lack of transparency of Meta's data collection practices raises questions about long-term data retention and user profiling. Additionally, the app's usage of WhatsApp during physical interactions, such as taking photos and recording videos in the network traffic analysis, stresses the questionable practices of data collection by Meta, prompting further concerns about how they might use or for how long they may retain the data. 


\subsection*{Trade-offs between Device Functionality and Privacy}
Meta has been continuously working to improve the privacy of its users by implementing measures to prevent unauthorized access to the glasses. As detailed in the device overview, one such measure introduced is to lock certain functionalities, such as having the AI read a message aloud, unless the paired phone is unlocked. Although this measurement improves user privacy, it diminishes the glasses' selling point of hands-free convenience. The requirement to unlock the phone before reading text messages aloud, for instance, negates the ideal hands-free interaction initially offered by the device, as the user would already have access to the message contents when unlocking the phone.

This trade-off represents a broader challenge where enhancing security and privacy controls comes at the cost of usability while optimizing ease of use can weaken protective measures. Even when some level of compromise may be needed, users would benefit more when these design decisions and their associated trade-offs are transparently communicated and available for users to decide. 


\subsection*{Policy and Ethical Considerations}
From a policy point of view, the observations from the network, Bluetooth, and app analyses show the lack of a clear regulatory framework and transparency surrounding the collection, storage, and sharing of data by wearable AR devices. Therefore, it is important to have clear guidelines on the retention periods of the data, the parties who have access to the data, and the specific purposes for processing the data. The encryption of the content prevents the interception of raw sensor data, addressing data confidentiality. However, the metadata analysis reveals substantial details about user activity and application usage, necessitating further technical solutions to address these exposures.

In addition, measures for bystander privacy remain underexplored, and there is no guarantee nor verification mechanism for bystander consent. Even though an LED notification is in place for bystanders on the Meta glasses, they can be easily missed or obstructed, as pointed out in the device overview. Users could bypass the security measures implemented in the glasses during video recording by blocking the LED indicator after starting a video recording, limiting bystander awareness. As a result, the glasses can be used as a discrete recording device, meaning that bystanders' data can still be easily and carelessly gathered, transmitted, and analyzed without their awareness or consent. This raises ethical questions about the transparency of AR devices in public spaces. Therefore, the risk from Meta's data collection, combined with the fact that bystanders, who have no direct relationship with the device or wearer, could be unknowingly captured in video and audio recordings, can lead to bystanders becoming part of a long-term analytic dataset.

The findings, emphasized by the device's heavy dependence on cloud services and the associated risks to bystanders, call for the need for stronger policies and strategies aimed at preserving the ability of users and bystanders to give informed consent. Device manufacturers and the wearers are both responsible for mitigating these concerns. 

